\subsection{Data Definition Language}
\begin{frame}[fragile]{INSERT INTO}
  \begin{block}{Adding Rows to a Table}
\begin{lstlisting}[style=sqlstyle]
-- Insert a single row (explicit column list)
INSERT INTO unit (id, name) VALUES ('T', 'ton');

-- Insert multiple rows at once
INSERT INTO unit (id, name)
VALUES ('KG', 'kilogram'),
       ('G',  'gram'),
       ('L',  'litre');

-- Insert from a SELECT result
INSERT INTO archive_orders
SELECT * FROM orders WHERE order_date < '2020-01-01';
\end{lstlisting}
  \end{block}

  \begin{itemize}
    \item Columns not listed in the column list receive their
          \texttt{DEFAULT} value, or \texttt{NULL} if no default is defined.
    \item \texttt{AUTO\_INCREMENT} columns (typically the primary key) are
          \textbf{omitted} from the column list -- the database assigns the
          next integer value automatically.
    \item When inserting into a table with \textbf{foreign key constraints},
          the referenced row must already exist -- otherwise the insert is
          rejected with a constraint violation error.
  \end{itemize}
\end{frame}

% ----------------------------------------------------------
\begin{frame}[fragile]{UPDATE}
  \begin{block}{Modifying Existing Rows}
\begin{lstlisting}[style=sqlstyle]
-- Update one column for a specific row
UPDATE employee
SET last_name = 'Miller'
WHERE id = 12;

-- Update multiple columns using expressions
UPDATE product
SET price         = price * 1.1,
    last_modified = NOW()
WHERE brand_id = 5;

-- DANGER: UPDATE without WHERE affects ALL rows!
UPDATE product SET price = 0; -- sets every product's price to zero!
\end{lstlisting}
  \end{block}

  \begin{alertblock}{Always Include a WHERE Clause}
    An \texttt{UPDATE} without \texttt{WHERE} modifies \textbf{every single
    row} in the table.  Before running any \texttt{UPDATE}, test your
    \texttt{WHERE} condition with a \texttt{SELECT} first to confirm you
    are targeting exactly the right rows.  Consider enabling
    \texttt{SET SQL\_SAFE\_UPDATES = 1} in MySQL to prevent accidental
    full-table updates.
  \end{alertblock}
\end{frame}

% ----------------------------------------------------------
\begin{frame}[fragile]{DELETE}
  \begin{block}{Removing Rows from a Table}
\begin{lstlisting}[style=sqlstyle]
-- Delete specific rows
DELETE FROM employee WHERE department_id = 5;

-- Delete with JOIN condition (MySQL syntax)
DELETE product
FROM product
INNER JOIN brand ON product.brand_id = brand.id
WHERE brand.country = 'USA';

-- DANGER: DELETE without WHERE empties the table!
DELETE FROM employee; -- removes every row!
\end{lstlisting}
  \end{block}

  \begin{itemize}
    \item \texttt{DELETE} removes rows but \textbf{keeps the table
          structure} (schema, indexes, constraints).  To remove the table
          entirely, use \texttt{DROP TABLE}.
    \item \texttt{TRUNCATE TABLE employee} is faster than
          \texttt{DELETE FROM employee} for emptying a table (it
          deallocates data pages directly), but it cannot be rolled back in
          some DBMSs and resets \texttt{AUTO\_INCREMENT} counters.
  \end{itemize}

  \begin{alertblock}{Always Include a WHERE Clause}
    The same warning as for \texttt{UPDATE} applies: test with a
    \texttt{SELECT} first.  Deleted rows cannot be recovered without a
    backup.
  \end{alertblock}
\end{frame}
